% =====================================================================
%  CV — Nicolás Plata Molano · format une colonne (inspiré de sa « Hoja de vida »)
%  Compilation : pdflatex CV_Nicolas_Plata_Molano_FR.tex
% =====================================================================
\documentclass[10pt,a4paper]{article}

\usepackage[utf8]{inputenc}
\usepackage[T1]{fontenc}
\usepackage[french]{babel}
\usepackage[sfdefault,book,lining]{FiraSans}
\usepackage[left=1.5cm,right=1.5cm,top=1.0cm,bottom=0.9cm]{geometry}
\usepackage{xcolor}
\usepackage{titlesec}
\usepackage{enumitem}
\usepackage{tabularx}
\usepackage[hidelinks]{hyperref}
\usepackage[expansion=false]{microtype}
\usepackage{amssymb}

\definecolor{navy}{HTML}{1F3F6E}
\definecolor{grey}{HTML}{5C6770}
\definecolor{rule}{HTML}{1F3F6E}

\pagestyle{empty}
\setlength{\parindent}{0pt}
\setlength{\parskip}{0pt}
\linespread{1.03}

\titleformat{\section}{\color{navy}\large\bfseries}{}{0pt}{}[\vspace{1pt}{\color{rule}\titlerule[0.6pt]}]
\titlespacing*{\section}{0pt}{6pt}{3pt}

% \entree{Titre — rôle}{Lieu / dates}{texte}
\newcommand{\entree}[3]{%
  \textbf{#1}\hfill{\small\itshape #2}\par\vspace{0.5pt}
  {\small #3}\par\vspace{3.5pt}}

\newcommand{\bati}{BaTiO\textsubscript{3}}

\begin{document}

% ======================= EN-TÊTE =======================
\begin{center}
  {\LARGE\bfseries NICOLÁS PLATA MOLANO}\\[2pt]
  {\color{grey}Élève-ingénieur en double diplôme — IMT Atlantique (ASCY) · Universidad Nacional de Colombia}\\[3pt]
  {\small \href{mailto:nicao511ng@gmail.com}{nicao511ng@gmail.com} \;|\; +33 7 45 57 71 66 \;|\; Nantes, France}\\[1pt]
  {\small \href{https://www.linkedin.com/in/nicolas-plata-molano}{linkedin.com/in/nicolas-plata-molano} \;|\;
  \href{https://github.com/NicolasPlata2004}{github.com/NicolasPlata2004} \;|\;
  \href{https://nicolasplata2004.github.io}{\textbf{Portfolio : nicolasplata2004.github.io}}}
\end{center}
\vspace{-4pt}
{\color{rule}\rule{\linewidth}{0.8pt}}
\vspace{2pt}

{\small\textbf{Recherche un stage de 2\textsuperscript{e} année} (4–5 mois dès avril 2027, ou 2–3 mois dès juin) en automatique,
robotique, systèmes embarqués ou automatisation industrielle — Nantes en priorité, mobile dans toute la France.}

% ======================= PROFIL =======================
\section{Profil}
{\small Élève-ingénieur en mécatronique, \textbf{5\textsuperscript{e} de ma promotion} (moyenne 4,4/5). J'identifie, commande et valide
des systèmes sur du matériel réel : quatre lois de commande testées sur banc, un bras robotique 2R, un objet IoT sur ESP32 et une ligne
pilotée par automate. En recherche (groupe KYMA), je travaille sur des capteurs piézoélectriques sans plomb (\bati).}

% ======================= FORMATION =======================
\section{Formation}
\entree{IMT Atlantique — Diplôme d'ingénieur généraliste (FISE 2A), double diplôme}{Nantes, 2026–2028}{%
Thématique ASCY : Automatique et Systèmes Cyber-physiques. Commande d'entreprise en cours : qualification d'une passerelle IoT industrielle
(MQTT, Modbus, OPC-UA, S7) pour une startup incubée à l'école.}
\entree{Universidad Nacional de Colombia — Ingénierie mécatronique}{Bogotá, 2021–2026}{%
\textit{Moyenne :} \textbf{4,4 / 5} \;|\; \textbf{5\textsuperscript{e} de la promotion} \;|\; \textit{Cursus validé :} 73,7 \%.
Moniteur académique ; formation en instrumentation industrielle, robotique et modélisation de systèmes dynamiques.}
\entree{I.E. Camilo Torres Restrepo — Baccalauréat}{Aguazul, Casanare}{%
Meilleure moyenne de l'établissement et score ICFES remarqué au niveau départemental.}

% ======================= EXPÉRIENCE =======================
\section{Recherche et projets}
\entree{Groupe de recherche KYMA (UNAL) — Chercheur, capteurs et commande}{2024–présent}{%
Projet \textit{Piezo} : synthèse et caractérisation du titanate de baryum (\bati) pour des capteurs piézoélectriques de force et de vibration
sans plomb. Conception électronique pour l'acquisition de signaux et l'IoT, algorithmes de commande robuste.}
\entree{Aéropendule — identification et quatre lois de commande}{Técnicas de Control, 2026}{%
Modèle non linéaire identifié et validé sous Simulink (R\textsuperscript{2} = 0,943) ; PID, H$\infty$, MRAC et mode glissant déployés sur Arduino et comparés sur le banc réel.}
\entree{Bras robotique 2R « ADANISS »}{Servomecanismos, 2026}{%
Trajectoire polaire paramétrique, cinématique inverse, identification des moteurs (R\textsuperscript{2} > 0,99) et PID discret sur ESP32 ;
carte électronique conçue et fabriquée par l'équipe.}
\entree{ThermaLink-C3 — supervision IoT de climatisation}{2026}{%
ESP32-C3 : détection de présence, température, clonage de trames IR et machine à états ; panneau web, \textbf{MQTT} et intégration
Home Assistant ; PCB sous KiCad. Économie estimée : environ 29 \%.}
\entree{Case Erector — automatisation industrielle}{Proyecto industrial, 2026}{%
Formeuse de caisses à trois axes (20 caisses/min) programmée en ladder (Studio 5000) sur automate émulé, validée par jumeau numérique (OPC DA, Siemens NX).}
\entree{UN Makerspace — Prototypage et IoT}{2023–présent}{%
Conception mécanique et prototypage par impression 3D ; automatisation de flux de travail (n8n) pour la traçabilité des projets.}
\entree{Contrôle thermique périopératoire — systèmes embarqués}{UNAL}{%
Boucle PI sur microcontrôleur PIC18F4550 et traitement numérique du signal en temps réel.}

% ======================= COMPÉTENCES =======================
\section{Compétences}
{\small
\textbf{Automatique :} identification, PID, H$\infty$, MRAC, mode glissant, MATLAB/Simulink (déploiement sur Arduino).\\
\textbf{Embarqué \& IoT :} ESP32, Arduino, PIC18F4550, C/C++, MQTT, Home Assistant, KiCad, Proteus, LTspice.\\
\textbf{Automatisme :} Studio 5000 (ladder), OPC, Modbus, jumeau numérique (Siemens NX), CADe SIMU, FluidSIM.\\
\textbf{CAO / FAO :} Fusion 360, Inventor, AutoCAD, Mastercam, impression 3D.\\
\textbf{Programmation :} Python (Pandas, NumPy, Matplotlib), MATLAB, C/C++, JavaScript/React, Git, LaTeX.\\
\textbf{Vulgarisation :} animations scientifiques avec Manim ; fondateur du blog « El Éter ».}

% ======================= LANGUES ET CERTIFICATIONS =======================
\section{Langues et certifications}
{\small
\textbf{Langues :} espagnol (langue maternelle) · anglais C1 (certification B2) · \textbf{français B2 (certification officielle)} · allemand A1.\\
\textbf{Certifications :} Diplôme en développement d'applications mobiles (MinTIC / Universidad Tecnológica de Pereira) ; fondamentaux de la
programmation et du développement logiciel (MinTIC) ; impression 3D métallique ; Python pour l'analyse de données.}

\end{document}
